\chapter{Полный физический модуль и вакуумный гессиан рангового поля}

\section{Решающий вопрос}

Первый гейт Тома VII допустил к проверке поле
\begin{equation}
 \Phi\in\Gamma\!\left(
 X,E_{\rm aff}\widehat\otimes\mathcal Y_{\rm phys}
 \right),
 \qquad
 E_{\rm aff}=\operatorname{Hom}(\mathbb C^4,\operatorname{im}P_3),
\end{equation}
с одним кривизностным действием. Минимальный плоский суррогат оказался
строго устойчивым. Теперь требуется восстановить фактический физический
модуль $H_{15}$ и проверить, создаёт ли ненулевой конечный оператор $D_F$
нецентральную отрицательную моду.

\section{Замороженный физический модуль}

Том V установил
\begin{equation}
 H_{15}=Q_L\oplus L_L\oplus u_R\oplus d_R\oplus e_R,
 \qquad
 15=6+2+3+3+1.
\end{equation}
На этом пакете существуют ровно три перенормируемых заряженных ребра
\begin{align}
 \lambda_u&\in\operatorname{Hom}_G(Q_L\otimes\widetilde H,u_R),\\
 \lambda_d&\in\operatorname{Hom}_G(Q_L\otimes H,d_R),\\
 \lambda_e&\in\operatorname{Hom}_G(L_L\otimes H,e_R).
\end{align}
Они попарно неэквивалентны как бимодули, поэтому
\begin{equation}
 \Lambda_{\rm ch}\simeq\mathbb C^3,
 \qquad
 \operatorname{End}_{\mathcal A_{\rm SM}-\mathcal A_{\rm SM}}
 (\Lambda_{\rm ch})\simeq\mathbb C^3.
 \label{eq:v7-h15-edge-commutant}
\end{equation}

Семейный модуль одноформ имеет комплексную размерность четыре, и
\begin{equation}
 \mathcal Y_{\rm phys}
 =\mathcal E_\rho\otimes\Lambda_{\rm ch},
 \qquad
 \dim_{\mathbb C}\mathcal Y_{\rm phys}=12.
\end{equation}
Следовательно, поле Тома VII до Real-удвоения имеет
\begin{equation}
 \dim_{\mathbb C}
 (E_{\rm aff}\otimes\mathcal Y_{\rm phys})
 =12\cdot12=144,
 \qquad
 \dim_{\mathbb R}=288.
 \label{eq:v7-full-rank-field-dimension}
\end{equation}

Один нормированный след существует, но не отождествляет три физических
ребра. После удаления общего направления остаются две вещественные
относительные амплитуды. Поэтому значения фонового $D_F$ не выводятся из
самой грассмановой связности.

\section{Радиальная вариация конечного оператора}

Ключевая проблема не требует знания численных юкавских коэффициентов.
Пусть $D_F\ne0$ является нечётной самосопряжённой частью базовой
суперсвязности и рассмотрим допустимую радиальную вариацию
\begin{equation}
 \widehat\eta=D_F.
\end{equation}
Она принадлежит тому же физическому модулю заряженных одноформ: запрещать
её означало бы заранее фиксировать длину поля и удалить именно направление,
которое должно менять ранг.

На постоянном пространственном фоне
\begin{equation}
 \mathbb A(t)=\mathbb A_{\rm sp}+(1+t)D_F.
\end{equation}
В секторе, где пространственная часть не создаёт смешанной кривизны,
кривизностная норма содержит
\begin{equation}
 \mathcal S_7(t)
 =\operatorname{tr}_{\rm norm}
 \left((1+t)^4D_F^4\right)+\text{постоянные члены}.
\end{equation}
Поэтому
\begin{equation}
 \left.\frac{d\mathcal S_7}{dt}\right|_{t=0}
 =4\operatorname{tr}_{\rm norm}(D_F^4)>0
 \qquad(D_F\ne0).
 \label{eq:v7-nonzero-df-nonstationarity}
\end{equation}

\begin{theorem}[Ненулевой фон не является нулевым вакуумом поля]
Если $D_F$ включён в $\mathbb A_0$, а поле $\Phi$ допускает физические
вариации в том же модуле одноформ, то конфигурация $\Phi=0$ не стационарна
для чистой нормы кривизны. Следовательно, вычисление её гессиана не является
тестом устойчивости вакуума: уже первая вариация ненулевая.
\end{theorem}

Знак производной не зависит от отношений $u,d,e$. Каждое ненулевое ребро
даёт положительный вклад. Свободные относительные амплитуды делают величину
градиента непредсказанной, но не устраняют саму нестационарность.

\section{Альтернативная ветвь: весь конечный оператор является полем}

Чтобы нулевой вакуум действительно соответствовал нулю, можно перенести
весь $D_F$ из фона в $\Phi$. Тогда базовая кривизна в постоянном конечном
секторе центральна или равна нулю, а квадратичная форма принимает вид
\begin{equation}
 \Hess_0\mathcal S_7(\eta,\eta)
 =2\|L_0\eta\|^2\ge0.
 \label{eq:v7-zero-df-positive-hessian}
\end{equation}
В минимальном аффинно-хиральном блоке это уже дало
\begin{equation}
 \Hess_0\mathcal S_7=\frac87I_{24}.
\end{equation}
После тензорного умножения на двенадцатимерный модуль
$\mathcal Y_{\rm phys}$ знак не меняется: в сырой нормировке полного следа
получаются двести восемьдесят восемь положительных направлений.

\begin{theorem}[Развилка стационарности и запуска]
Для родителя
$\mathcal S_7=\|\mathbb A_\Phi^2\|^2$ существуют две естественные
типизации нулевого сектора.
\begin{enumerate}
 \item При фиксированном ненулевом $D_F$ точка $\Phi=0$ не стационарна.
 \item При включении $D_F$ в динамическое поле точка $\Phi=0$ стационарна,
 но её физический гессиан неотрицателен и в минимальном полном модуле строго
 положителен.
\end{enumerate}
Ни одна ветвь не даёт эндогенную отрицательную моду нулевого вакуума.
\end{theorem}

\section{Почему очевидные исправления запрещены}

Можно формально сделать ненулевой $D_F$ стационарным, заменив действие на
\begin{equation}
 \|\mathbb F_\Phi-\mathbb F_{\rm ref}\|^2
\end{equation}
или добавив отрицательный квадратичный член. Но $\mathbb F_{\rm ref}$,
массовый знак и его нормировка не следуют из принятого кандидата. Их выбор
после обнаружения положительного гессиана нарушает $P1$ и повторяет ручной
quench Тома VI.

Ограничение вариаций условием
\begin{equation}
 \langle D_F,\eta\rangle=0
\end{equation}
также не является решением: оно фиксирует радиальную норму и удаляет
направление смены ранга до динамики.

\begin{nogocriterion}[Запрет фоновой подстановки]
Нельзя одновременно считать $D_F$ фиксированным ненулевым фоном, называть
$\Phi=0$ стационарным пустым вакуумом и разрешать $\Phi$ пробегать полный
физический модуль одноформ. Эти три требования несовместимы для чистой нормы
кривизны.
\end{nogocriterion}

\section{Вердикт по первому родителю Тома VII}

\begin{center}
\begin{tabular}{p{0.12\textwidth}p{0.18\textwidth}p{0.58\textwidth}}
\toprule
Пункт & Статус & Причина \\
\midrule
$P0$ & пройден типизационно & общий модуль существует и имеет размерность
$144$ над $\mathbb C$ до Real-завершения \\
$P1$ & понижен & один след существует, но ненулевой $D_F$ содержит две
нефиксированные относительные амплитуды рёбер \\
$P2$ & провален & ненулевой фон нестационарен, нулевой фон имеет
неотрицательный гессиан \\
$P3$--$P5$ & остановлены & переход к endpoint, скорости и масштабу запрещён \\
$P6$ & пройден & отрицательный исход воспроизводим и имеет стоп-критерий \\
\bottomrule
\end{tabular}
\end{center}

\begin{theorem}[Закрытие чистой кривизностной ветви]
Первый родитель Тома VII в форме чистой нормы кривизны не порождает
физическую ранговую неустойчивость нулевого сектора. При ненулевом $D_F$
нет стационарного нулевого вакуума; при нулевом $D_F$ нет отрицательной
моды. Ветвь закрывается до поиска локализованного объекта.
\end{theorem}

Этот результат не закрывает весь Том VII. Он устанавливает требование к
следующему кандидату: стационарный ненулевой фон и его потенциальная
неустойчивость должны следовать из нового функционала до вычисления знака,
а не из вычитания фоновой кривизны после результата.

\begin{protocol}
\begin{enumerate}
 \item комплексная размерность $E_{\rm aff}$: \textbf{$12$};
 \item комплексная размерность $\mathcal Y_{\rm phys}$: \textbf{$12$};
 \item полный комплексный модуль поля: \textbf{$144$};
 \item полный вещественный касательный модуль: \textbf{$288$};
 \item число неэквивалентных заряженных рёбер: \textbf{$3$};
 \item остаточных относительных вещественных направлений: \textbf{$2$};
 \item ненулевой $D_F$ и стационарность $\Phi=0$: \textbf{несовместимы};
 \item нулевой $D_F$ и отрицательная мода: \textbf{нет};
 \item чистая кривизностная ветвь: \textbf{закрыта отрицательно};
 \item следующий допустимый шаг: \textbf{новый заранее заданный функционал,
 имеющий стационарный ненулевой фон без ручного вычитания}.
\end{enumerate}
\end{protocol}

Машинный сертификат сохраняется в
\path{s2t_v7_full_physical_rank_field_hessian_gate_results.json}.